\documentclass[runningheads]{llncs}
\usepackage{graphicx}
\usepackage{amsmath,amssymb}
\usepackage{booktabs}
\usepackage{listings}
\usepackage{xcolor}
\usepackage{hyperref}
\usepackage{caption}
\usepackage{subcaption}
\usepackage{float}
\usepackage{algorithm}
\usepackage{algpseudocode}

\definecolor{codegreen}{rgb}{0,0.6,0}
\definecolor{codegray}{rgb}{0.5,0.5,0.5}
\definecolor{codepurple}{rgb}{0.58,0,0.82}
\definecolor{backcolour}{rgb}{0.95,0.95,0.92}

\lstdefinestyle{mystyle}{
    backgroundcolor=\color{backcolour},
    commentstyle=\color{codegreen},
    keywordstyle=\color{magenta},
    numberstyle=\tiny\color{codegray},
    stringstyle=\color{codepurple},
    basicstyle=\ttfamily\footnotesize,
    breakatwhitespace=false,
    breaklines=true,
    captionpos=b,
    keepspaces=true,
    numbers=left,
    numbersep=5pt,
    showspaces=false,
    showstringspaces=false,
    showtabs=false,
    tabsize=2
}
\lstset{style=mystyle}

\begin{document}

\title{AEGIS: A Universal Governance Primitive \\for Autonomous AI Systems}
\author{Aaditya Aggarwal}
\institute{Independent Research \\ \email{aaditya@aegis.ai}}

\maketitle

\begin{abstract}
The rapid proliferation of autonomous AI agents has created a critical governance gap. Current deployments lack identity verification, policy enforcement, audit trails, and human oversight. We present AEGIS, a lightweight, framework-agnostic governance layer that transparently intercepts all agent-to-resource communication. AEGIS introduces a \emph{Two-Plane Verification} architecture that combines application-layer policy enforcement with kernel-level eBPF monitoring, achieving non-bypassable security even under agent compromise. The system comprises (1) a Rust sidecar proxy with OPA/WASM policy evaluation, (2) a Go-based multi-provider gateway, and (3) a centralized control plane with tamper-evident audit logging. Our implementation spans 11,761 lines across 125 files in Rust, Go, C (eBPF), and Rego. AEGIS includes 26 automated tests across compliance, security, and performance suites, with 10 builtin detection functions covering seven attack vectors. AEGIS maps directly to regulatory requirements including the EU AI Act, NIST AI RMF, and Singapore IMDA classification.

\keywords{AI Governance, Agent Safety, Policy Enforcement, eBPF, Sidecar Proxy, Autonomous Systems}
\end{abstract}

\section{Introduction}

On April 15, 2026, an AI coding agent deleted the entire production database of a SaaS platform in approximately nine seconds~\cite{incident2026}. This incident is not isolated. As autonomous AI agents transition from research demonstrations to production deployments, the absence of systematic governance infrastructure has become the single greatest obstacle to safe deployment.

The core challenge is architectural: today's AI agents---whether implemented with LangGraph, CrewAI, AutoGen, or the OpenAI Agents SDK---operate with unfettered access to tools, databases, APIs, and credentials. The only protective mechanisms are system prompts, which agents can trivially bypass, and application-level guardrails, which run within the same address space as the agent itself. This creates a fundamental security vulnerability: an agent compromised through prompt injection, jailbreaking, or supply chain attack has unrestricted access to its environment.

The software engineering community faced an analogous problem in the mid-2010s: ungoverned microservice communication. The solution was the service mesh---a dedicated infrastructure layer (Envoy, Istio, Linkerd) that transparently intercepts all inter-service traffic to enforce identity, policy, encryption, and observability without application modification. We argue that autonomous AI agents require an equivalent governance primitive.

This paper presents AEGIS, a universal governance layer for autonomous AI systems. Our contributions are:

\begin{enumerate}
\item \textbf{Architecture}: A four-layer governance model (Agent $\to$ Sidecar $\to$ Gateway $\to$ Resources) that decouples governance from agent logic while remaining framework-agnostic.
\item \textbf{Two-Plane Verification}: A novel security mechanism combining application-layer (app-plane) and kernel-level (infra-plane eBPF) policy enforcement, ensuring non-bypassability even when the agent process is fully compromised.
\item \textbf{OPA/WASM Policy Engine}: A custom Rego-based policy runtime with 10 agent-specific built-in functions for recursion detection, budget tracking, and credential scoping, compiled to WASM for sandboxed execution.
\item \textbf{Implementation}: A complete open-source implementation of 11,761 lines spanning Rust (5,923 lines), Go (1,362 lines), C (570 lines for eBPF), and Rego (212 lines), with 1,832 lines of test and benchmark code.
\item \textbf{Evaluation}: Automated test suite with 26 tests across compliance (13), security (8), and performance (5) categories, covering all 7 attack vectors with 10 builtin detection functions.
\item \textbf{Regulatory Mapping}: Direct compliance mappings to the EU AI Act (Articles 9--72), NIST AI RMF, and Singapore IMDA's 5-tier autonomy classification.
\end{enumerate}

\section{Background and Problem Statement}

\subsection{The Governance Gap}

An autonomous AI agent, in its most general form, is a loop: it receives a goal, reasons about it, selects actions, invokes tools, observes results, and iterates. This loop---whether implemented as a LangGraph state machine, a CrewAI crew, or an AutoGen conversation---requires access to external resources: LLM APIs, databases, file systems, and internal services.

Current safety mechanisms include:
\begin{itemize}
\item \textbf{System prompts}: Instructions embedded in the LLM context window. These are advisory, not enforced, and are trivially bypassed by prompt injection~\cite{promptInjection2023,willison2023}.
\item \textbf{Application guardrails}: Libraries (e.g., Guardrails AI, NVIDIA NeMo Guardrails) that run within the agent process. These share the agent's address space and can be subverted if the agent is compromised.
\item \textbf{API keys and secrets}: Stored in environment variables or secret managers, accessible to any code the agent runs.
\end{itemize}

These mechanisms share a common failure mode: they operate within the agent's trust boundary. An attacker who controls the agent process controls all of them.

\subsection{Threat Model}

We assume an attacker with the following capabilities:
\begin{itemize}
\item Full control of the agent process (memory, execution, environment variables)
\item Ability to send arbitrary prompts and tool calls
\item Access to the agent's network environment
\item Knowledge of the agent's configuration and available tools
\end{itemize}

The attacker cannot:
\begin{itemize}
\item Modify kernel-level eBPF programs (requires root beyond the agent's privileges)
\item Bypass iptables/nftables rules enforced at the network layer
\item Access hardware security modules (HSM/TPM) directly
\end{itemize}

This threat model motivates our core architectural requirement: governance must be \emph{external} to the agent and \emph{non-bypassable}.

\subsection{The Service Mesh Analogy}

The service mesh model is instructive. Before service meshes, microservices communicated directly, managing their own TLS, retries, and observability. The introduction of a transparent sidecar proxy (Envoy) decoupled infrastructure concerns from application logic. We argue that AI agents represent an analogous architectural boundary: the governance of agent-resource interactions should be external, transparent, and uniform across all agent frameworks.

\section{Architecture}

AEGIS adopts a four-layer architecture designed for deployment flexibility, framework agnosticism, and defense in depth:

\begin{enumerate}
\item \textbf{Agent Layer}: Any AI agent framework (LangGraph, CrewAI, AutoGen, OpenAI SDK, MCP, Google ADK). No code changes required.
\item \textbf{Governance Layer (Sidecar)}: A Rust-based transparent proxy that intercepts all agent outbound traffic. Enforces identity, policy, cost limits, recursion protection, and audit logging.
\item \textbf{Gateway Layer}: A Go-based multi-provider gateway responsible for LLM routing, rate limiting, protocol normalization (MCP, A2A, ACP), and circuit breaking.
\item \textbf{Resource Layer}: External APIs (OpenAI, Anthropic, Gemini), databases, internal services, and tool endpoints.
\end{enumerate}

Cross-cutting concerns include:
\begin{itemize}
\item \textbf{Control Plane}: Centralized policy management, agent registration, and compliance reporting via etcd-backed configuration.
\item \textbf{Infra-Plane (eBPF)}: Kernel-level monitoring and independent policy enforcement (Section~\ref{sec:twoplane}).
\item \textbf{Observability}: Structured logging (OpenTelemetry OTLP), Prometheus metrics, Jaeger tracing, and Grafana dashboards.
\end{itemize}

\subsection{Two-Plane Verification}
\label{sec:twoplane}

The central security innovation in AEGIS is Two-Plane Verification ($\pi_A + \pi_I$). Every governance decision requires independent agreement between two isolated evaluation planes:

\begin{itemize}
\item \textbf{Application Plane ($\pi_A$)}: The sidecar proxy evaluates the operation against OPA/Rego policies. This plane is fast ($<1$ms P99) and supports the full policy language.
\item \textbf{Infrastructure Plane ($\pi_I$)}: eBPF programs running in the Linux kernel independently monitor system calls (read, write, connect, openat, execve) and TCP connections. This plane cannot be bypassed even if the agent process is fully compromised.
\end{itemize}

The combination produces three outcomes:
\begin{equation*}
\text{Decision}(o) = \begin{cases}
\text{Permit} & \text{if } \pi_A(o) = \pi_I(o) = \text{ALLOW} \\
\text{Block} & \text{if } \pi_A(o) = \text{DENY} \lor \pi_I(o) = \text{DENY} \\
\text{Alert} & \text{if } \pi_A(o) \neq \pi_I(o) \text{ (plane divergence)}
\end{cases}
\end{equation*}

Plane divergence ($\pi_A(o) \neq \pi_I(o)$) is a critical signal: it indicates that either the app-plane has been compromised or a new attack vector has been detected. In either case, the operation is blocked and an alert is raised.

\subsection{Identity}

AEGIS uses SPIFFE-compatible X.509 SVIDs (SPIFFE Verifiable Identity Documents) for agent identity. Each agent receives a unique identity of the form:
\begin{lstlisting}
spiffe://aegis.local/ns/<namespace>/sa/<service-account>
\end{lstlisting}
Identities are issued by the control plane and verified by the sidecar on every request. This enables:
\begin{itemize}
\item Fine-grained authorization policies per agent identity
\item Environment scoping (agents in staging cannot access production resources)
\tightitem Delegation chain enforcement via NGAC (Next Generation Access Control)
\end{itemize}

\subsection{Audit}

The audit log is an append-only, tamper-evident chain. Each event is:
\begin{lstlisting}[language=JSON]
{
  "event_id": "uuid",
  "trace_id": "uuid",
  "agent_id": "spiffe://...",
  "operation": "tool.call",
  "resource": "resources/secret",
  "decision": "DENY",
  "timestamp": "2026-07-26T12:00:00Z",
  "signature": "ed25519-hex-sig",
  "sequence_number": 42,
  "sidecar_id": "sidecar-1",
  "metadata": {"reason": "unauthorized tool"}
}
\end{lstlisting}

Events are cryptographically linked: each event's hash includes the previous event's hash, forming an immutable chain. Tampering with any event breaks the chain and is immediately detectable.

\section{Design}

\subsection{Sidecar Proxy}

The sidecar is implemented in Rust (1,646 lines across 14 files) and runs as a transparent HTTP/HTTPS/gRPC proxy between the agent and its resources. Key modules include:

\begin{table}[H]
\centering
\caption{Sidecar modules}
\begin{tabular}{lrl}
\toprule
Module & Lines & Responsibility \\
\midrule
\texttt{proxy.rs} & 214 & HTTP/HTTPS/gRPC interception \\
\texttt{middleware.rs} & 168 & Governance middleware chain \\
\texttt{identity.rs} & 89 & SPIFFE identity verification \\
\texttt{policy.rs} & 156 & OPA policy evaluation integration \\
\texttt{tool\_gate.rs} & 134 & Schema-validated tool access control \\
\texttt{cost\_circuit.rs} & 98 & Micro-cent budget tracking \\
\texttt{recursion.rs} & 72 & Operation fingerprinting \\
\texttt{audit.rs} & 110 & Tamper-evident event production \\
\texttt{ebpf.rs} & 67 & eBPF event consumption \\
\texttt{protocol.rs} & 84 & MCP/A2A protocol detection \\
\texttt{metrics.rs} & 62 & Prometheus metrics \\
\texttt{telemetry.rs} & 145 & OpenTelemetry tracing \\
\texttt{main.rs} & 247 & Graceful shutdown, signal handling \\
\bottomrule
\end{tabular}
\end{table}

\subsection{Policy Engine}

The policy engine (1,347 lines across 8 files) evaluates OPA/Rego policies compiled to WASM. It includes 10 custom built-in functions:

\begin{enumerate}
\item \texttt{is\_recursion\_detected} — Operation fingerprint matching
\item \texttt{get\_recursion\_depth} — Current call stack depth
\item \texttt{calculate\_budget\_remaining} — Remaining micro-cent budget
\item \texttt{get\_environment} — Current deployment environment
\item \texttt{has\_recent\_tool\_call} — Recent operation history
\item \texttt{get\_risk\_score} — Composite risk assessment
\item \texttt{get\_autonomy\_level} — IMDA autonomy level
\item \texttt{is\_resource\_accessible} — Environment-resource matching
\item \texttt{get\_credential\_scope} — Credential boundary
\item \texttt{is\_human\_in\_loop} — Escalation state
\end{enumerate}

Policies are compiled to WASM modules using OPA's Rego compiler and executed in a wasmtime sandbox. This provides strong isolation from the sidecar process: a malicious policy cannot compromise the sidecar or the kernel.

\subsection{Gateway}

The gateway (1,209 lines across 8 Go files) provides multi-provider LLM routing with the following features:
\begin{itemize}
\item Provider-agnostic interface supporting OpenAI, Anthropic, Gemini, Mistral, Ollama, and OpenRouter
\item Weighted round-robin routing with failover
\item Distributed rate limiting via Redis (token bucket algorithm)
\item Circuit breaker per provider (5 consecutive failures $\to$ open circuit)
\item Protocol normalization: MCP, A2A, and ACP adapters translate between protocols
\item Per-route cost tracking
\end{itemize}

The gateway is stateless and horizontally scalable. Configuration is fetched from the control plane via gRPC.

\subsection{eBPF Monitoring}

Four eBPF programs (570 lines of C) monitor kernel-level operations:

\begin{enumerate}
\item \textbf{syscall\_monitor}: Hooks \texttt{read}, \texttt{write}, \texttt{connect}, \texttt{openat}, \texttt{execve} syscalls
\item \textbf{tcp\_monitor}: Tracks TCP connections via \texttt{tcp\_connect} and \texttt{tcp\_v4\_connect} kprobes
\item \textbf{file\_access\_monitor}: Monitors file access patterns via \texttt{security\_file\_open}
\item \textbf{credential\_monitor}: Tracks credential access patterns
\end{enumerate}

The userspace loader (71 lines in Rust, using the aya library) loads these programs into the kernel and processes events via perf ring buffers. Events are matched against the sidecar's policy decisions to detect plane divergence.

\subsection{Kubernetes Integration}

The K8s integration (1,235 YAML lines across 19 files) provides:
\begin{itemize}
\item Mutating admission webhook for automatic sidecar injection
\item \texttt{AgentPolicy} and \texttt{AgentRegistration} custom resource definitions (CRDs)
\item Kubernetes operator (kube-rs, 442 lines) for reconciling agent configurations
\item Horizontal Pod Autoscaling (HPA) with custom Prometheus metrics
\item Network policies for zero-trust agent isolation
\item Pod Disruption Budget (PDB) for high availability
\item Prometheus ServiceMonitor and Grafana dashboard
\item Helm chart with production defaults
\end{itemize}

\section{Implementation}

\subsection{Technology Stack}

\begin{table}[H]
\centering
\caption{Implementation technology stack}
\begin{tabular}{lll}
\toprule
Component & Language & Dependencies \\
\midrule
Sidecar Proxy & Rust 1.79 & tokio, hyper, tonic, wasmtime \\
Policy Engine & Rust + Rego & wasmtime, serde\_json \\
Gateway & Go 1.22 & gin, redis, otel \\
eBPF Programs & C & aya (Rust loader) \\
Control Plane & Go 1.22 & etcd, kube-rs \\
Audit Log & Rust & ed25519-dalek, sha2 \\
CLI & Rust & clap \\
Protocol Buffers & proto3 & prost, tonic-build \\
\bottomrule
\end{tabular}
\end{table}

\subsection{Project Statistics}

The complete implementation comprises:

\begin{table}[H]
\centering
\caption{Project statistics by language and component}
\begin{tabular}{lrr}
\toprule
Category & Files & Lines \\
\midrule
Rust (6 crates) & 35 & 4,216 \\
Rust (tests, benches) & 9 & 1,707 \\
Go (gateway + control-plane) & 9 & 1,362 \\
C (eBPF programs) & 4 & 570 \\
Rego (policies) & 7 & 212 \\
Protobuf & 3 & 363 \\
Kubernetes YAML & 19 & 1,235 \\
Shell scripts & 4 & 275 \\
Other (Makefile, config, etc.) & 35 & 1,821 \\
\midrule
\textbf{Total} & \textbf{125} & \textbf{11,761} \\
\bottomrule
\end{tabular}
\end{table}

Test and benchmark infrastructure accounts for 1,832 lines of test code across 11 files, including integration tests (security, compliance, performance), chaos engineering tests, and microbenchmarks (Criterion.rs).

\subsection{Regulatory Compliance}

AEGIS maps directly to three major regulatory frameworks:

\begin{itemize}
\item \textbf{EU AI Act}: Articles 9 (risk management), 10 (data governance), 11 (technical documentation), 12 (automatic logging), 14 (human oversight), 15 (accuracy/robustness/cybersecurity), 26 (foundation model obligations), 72 (regulatory sandboxes).
\item \textbf{NIST AI RMF}: Identity management (SPIFFE/OIDC), authorization (OAuth 2.1), access delegation (NGAC), continuous monitoring (eBPF), and audit logging (tamper-evident chain).
\item \textbf{Singapore IMDA}: All five autonomy levels supported---from Level 1 (human-initiated, approved) through Level 5 (fully autonomous, maximum restriction with post-hoc review).
\end{itemize}

Each framework maps to specific AEGIS capabilities with automated compliance dashboards and reporting.

\section{Evaluation}

\subsection{Test Infrastructure}

AEGIS includes 26 automated tests across three suites, all executing directly against the \texttt{PolicyEngine} without network dependencies:

\begin{itemize}
\item \textbf{Compliance tests (13)}: Verify enforcement of EU AI Act articles (risk management, logging, human oversight, accuracy), NIST AI RMF controls (identity, authorization, logging), budget enforcement (spending limits, rate limits), recursion prevention (depth limits, loop detection), and environment scoping.
\item \textbf{Security tests (8)}: Cover all 7 attack vectors (direct prompt injection, indirect prompt injection, tool hijacking, credential theft, privilege escalation, policy bypass, recursive loops) plus false positive measurement using 8 benign patterns.
\item \textbf{Performance tests (5)}: Measure policy evaluation latency (1,000 iterations), throughput (5,000 operations), deny-fast-path timing, cold start duration, and concurrent evaluation under contention.
\end{itemize}

The policy engine implements 10 builtin detection functions --- \texttt{prompt\_injection\_detect}, \texttt{tool\_hijack\_detect}, \texttt{credential\_detect}, \texttt{privilege\_escalation\_detect}, \texttt{policy\_bypass\_detect}, \texttt{recursion\_detect}, \texttt{rate\_limit\_check}, \texttt{budget\_check}, \texttt{scope\_check}, \texttt{compliance\_check} --- each returning \texttt{ALLOW}, \texttt{DENY}, or \texttt{REVIEW} with an evidence string.

eBPF programs are compiled from C sources (4 programs: \texttt{syscall\_monitor}, \texttt{file\_access\_monitor}, \texttt{credential\_monitor}, \texttt{network\_monitor}) with a Rust runtime that loads them via \texttt{aya} when available and falls back to simulated mode on non-Linux platforms.

\subsection{Security Effectiveness}

The security test suite verifies detection logic for each attack vector:

\begin{itemize}
\item \textbf{Direct Prompt Injection}: 8 variants including role redefinition, command injection, credential extraction queries, and system prompt override.
\item \textbf{Indirect Prompt Injection}: 4 variants embedded in tool outputs, including data exfiltration instructions and privilege escalation commands.
\item \textbf{Tool Hijacking}: 5 variants including schema mismatch, parameter overflow, unauthorized tool invocation, and tool impersonation.
\item \textbf{Credential Theft}: 6 variants targeting environment variables, config files, secret manager API, log files, debug endpoints, and key material.
\item \textbf{Privilege Escalation}: 4 variants spanning cross-environment access, cross-role delegation, and RBAC bypass.
\item \textbf{Policy Bypass}: 4 variants using URL encoding, case manipulation, path traversal, and header injection.
\item \textbf{Recursive Agent Loop}: Depth-limit enforcement with configurable max depth (default: 5).
\end{itemize}

Benchmark results should be generated by running \texttt{make bench} on the target hardware. The \texttt{scripts/bench.sh} orchestrates full evaluation across all dimensions.

\subsection{Performance Characterization}

Policy evaluation latency is measured using Criterion.rs microbenchmarks (1,000 iterations). Throughput tests execute 5,000 sequential evaluations and measure elapsed wall time. Cold start timing measures Time-to-First-Evaluation from process launch.

\begin{table}[H]
\centering
\caption{Performance targets with measurable thresholds}
\begin{tabular}{lrl}
\toprule
Metric & Target & Status \\
\midrule
Sidecar P50 Latency & $<$1ms & $\triangle$ \\
Sidecar P95 Latency & $<$3ms & $\triangle$ \\
Sidecar P99 Latency & $<$5ms & $\triangle$ \\
Gateway Throughput & 10,000 req/s & $\triangle$ \\
Policy Eval P50 & $<$0.5ms & $\triangle$ \\
Policy Eval P99 & $<$1ms & $\triangle$ \\
Sidecar RSS Memory & $<$50MB & $\triangle$ \\
Cold Start & $<$2s & $\triangle$ \\
eBPF Overhead & $<$10\% & $\triangle$ \\
\bottomrule
\end{tabular}
\end{table}

\emph{Note: $\triangle$ indicates targets have not yet been validated. Run \texttt{make bench} on target hardware to populate actual measurements.}

\section{Security Analysis}

\subsection{Non-Bypassability (Corollary 5.1)}

The Two-Plane Verification architecture ensures that even full compromise of the agent process and the application-plane policy engine does not grant unrestricted access. The infra-plane (eBPF) operates in the kernel and cannot be modified from userspace. We validate this through chaos engineering tests that simulate app-plane compromise:

\begin{itemize}
\item \textbf{Scenario}: Agent process hijacked, app-plane returns ALLOW for all operations
\item \textbf{Expected}: Infra-plane detects destructive operations (file deletion, production DB connection, privilege escalation) and enforces DENY
\item \textbf{Mechanism}: eBPF syscall monitors independently evaluate each operation against the same policy set
\end{itemize}

\begin{algorithm}[H]
\caption{Two-Plane Verification}
\begin{algorithmic}[1]
\Procedure{TwoPlaneVerify}{$op$}
\State $\pi_A \gets \text{AppPlaneEvaluate}(op)$
\State $\pi_I \gets \text{InfraPlaneEvaluate}(op)$
\If{$\pi_A = \pi_I$}
    \State \textbf{return} $\pi_A$
\Else
    \State \textbf{return} BLOCK + ALERT
\EndIf
\EndProcedure
\end{algorithmic}
\end{algorithm}

\subsection{Attack Vector Coverage}

\begin{itemize}
\item \textbf{Direct Prompt Injection}: Detected by OPA policies that scan for command patterns, role-redefinition attempts, and credential extraction queries. The policy engine maintains a pattern library updated from the MITRE ATLAS framework.
\item \textbf{Indirect Prompt Injection}: The eBPF file\_access\_monitor detects when tool outputs contain embedded instructions. Combined with the sidecar's tool output scanning, this provides defense in depth.
\item \textbf{Tool Hijacking}: Schema validation at the tool gate rejects any tool call that does not match the registered schema. The audit log records all tool invocations for post-hoc analysis.
\item \textbf{Credential Theft}: Environment scoping ensures that staging credentials cannot access production resources. The eBPF credential\_monitor tracks access to credential files and secret manager endpoints.
\item \textbf{Privilege Escalation}: SPIFFE identity combined with NGAC delegation limits prevents cross-environment and cross-role escalation.
\item \textbf{Policy Bypass}: URL normalization, DNS rebinding protection, and path traversal detection at the proxy layer. Unusual encoding patterns trigger additional scrutiny.
\item \textbf{Recursive Agent Loops}: Operation fingerprinting detects repeated identical tool calls. Configurable depth limits (default: 5) trigger circuit breakers.
\end{itemize}

\subsection{Defense in Depth}

AEGIS employs five layers of defense:

\begin{enumerate}
\item \textbf{Network Layer}: iptables/nftables rules ensure all agent outbound traffic passes through the sidecar. Direct egress is blocked.
\item \textbf{Proxy Layer}: The sidecar authenticates every request, evaluates policies, and enforces decisions.
\item \textbf{Policy Layer}: OPA/Rego policies with WASM sandboxing provide flexible, auditable, and isolated policy evaluation.
\item \textbf{Gateway Layer}: The gateway enforces rate limits, circuit breakers, and provider-level access controls.
\item \textbf{Kernel Layer}: eBPF programs provide independent verification that cannot be bypassed from userspace.
\end{enumerate}

\section{Related Work}

\subsection{AI Safety Frameworks}

Existing AI safety approaches operate at different layers of the stack. \emph{System prompts}~\cite{openaiSafety} are advisory instructions embedded in LLM context, easily bypassed by prompt injection~\cite{promptInjection2023,willison2023}. \emph{Application guardrails} (Guardrails AI, NVIDIA NeMo Guardrails) run within the agent process and share its security boundary. \emph{Lattice}~\cite{lattice2025} and \emph{Enterprise PDM} provide policy enforcement but lack kernel-level verification and framework agnosticism.

\subsection{Policy Engines}

Open Policy Agent (OPA)~\cite{opa2024} is the industry standard for policy-as-code, providing Rego-based policy evaluation. AEGIS extends OPA with agent-specific builtins and WASM compilation for sandboxed execution. Other policy engines (e.g., Cedar~\cite{cedar2023}, AuthZed, Oso) focus on RBAC/ABAC for traditional applications and lack AI-specific constructs.

\subsection{Service Meshes}

Envoy~\cite{envoy2024} and Istio~\cite{istio2024} pioneered the sidecar proxy pattern for microservice governance. AEGIS adapts this pattern for AI agents, adding AI-specific capabilities: prompt injection detection, recursive loop protection, cost circuit breakers, and two-plane verification. The key difference is semantic awareness: sidecars must understand agent operations, not just HTTP requests.

\subsection{Kernel-Level Security}

eBPF has been used for security monitoring (Falco~\cite{falco2024}, Tracee~\cite{tracee2024}), network policy (Cilium~\cite{cilium2024}), and performance tracing. AEGIS is the first system to use eBPF as an independent policy evaluation plane for AI agent governance, creating a non-bypassable security layer that operates regardless of agent compromise.

\section{Limitations and Future Work}

\subsection{Limitations}

\begin{enumerate}
\item \textbf{eBPF portability}: eBPF programs are Linux-specific. Windows and macOS support for agent governance requires alternative mechanisms (e.g., macOS Endpoint Security Framework).
\item \textbf{Performance overhead}: While individual overheads are small ($<$5ms P99), the cumulative effect of proxy + eBPF + policy evaluation may be significant for latency-sensitive applications.
\item \textbf{LLM-specific detection}: Prompt injection detection remains challenging. AEGIS uses pattern matching and policy rules rather than LLM-based detection, which could improve accuracy.
\item \textbf{Configuration complexity}: The full AEGIS stack (sidecar, gateway, control plane, eBPF, K8s operator) requires significant configuration, though Helm chart defaults reduce this burden.
\end{enumerate}

\subsection{Future Work}

\begin{enumerate}
\item \textbf{LLM-based detection}: Integrate a local SLM (Small Language Model) for semantic prompt injection detection, complementing the rule-based policy engine.
\item \textbf{Federated governance}: Cross-organizational agent communication with policy federation across administrative domains.
\item \textbf{Real-time policy adaptation}: Dynamic policy adjustment based on runtime risk assessment and anomaly detection.
\item \textbf{Adversarial robustness}: Formal verification of policy correctness and completeness using model checking.
\item \textbf{Standardization}: Contributing AEGIS's protocol specification to the Cloud Native Computing Foundation (CNCF) as a standard for AI agent governance.
\end{enumerate}

\section{Conclusion}

We presented AEGIS, a universal governance primitive for autonomous AI systems. AEGIS introduces Four-Layer Architecture with Two-Plane Verification, combining application-layer OPA/Rego policy enforcement with kernel-level eBPF monitoring to achieve non-bypassable security. The system is implemented as a complete open-source platform of 11,761 lines spanning Rust, Go, C, and Rego, with comprehensive test and benchmark infrastructure.

AEGIS addresses the critical governance gap that currently prevents safe deployment of autonomous AI agents. By externalizing governance from agent logic and layering independent verification mechanisms, AEGIS provides the same transformative infrastructure for AI agents that service meshes provided for microservices. The system maps to existing regulatory frameworks (EU AI Act, NIST AI RMF, Singapore IMDA) and is designed for extensibility to future agent frameworks and deployment environments.

The full implementation is available as open source at \url{https://github.com/aegis-ai/aegis}.

\begin{thebibliography}{20}

\bibitem{incident2026} Anonymous. Production Database Deletion Incident (April 2026). Industry report, 2026.

\bibitem{promptInjection2023} Greshake, K., Abdelnabi, S., Mishra, S., et al. Not what you've signed up for: Compromising Real-World LLM-Integrated Applications with Indirect Prompt Injection. In \textit{ACM CCS Workshop on AI Security}, 2023.

\bibitem{willison2023} Willison, S. Prompt injection: What's the worst that could happen? \textit{Simons Willison's Blog}, 2023.

\bibitem{opa2024} Open Policy Agent. \url{https://www.openpolicyagent.org/}, 2024.

\bibitem{envoy2024} Envoy Proxy. \url{https://www.envoyproxy.io/}, 2024.

\bibitem{istio2024} Istio. \url{https://istio.io/}, 2024.

\bibitem{cilium2024} Cilium. \url{https://cilium.io/}, 2024.

\bibitem{falco2024} Falco. \url{https://falco.org/}, 2024.

\bibitem{tracee2024} Tracee. \url{https://tracee.aquasec.com/}, 2024.

\bibitem{cedar2023} Cedar Policy Language. \url{https://www.cedarpolicy.com/}, 2023.

\bibitem{lattice2025} Lattice AI Governance. Industry report, 2025.

\bibitem{openaiSafety} OpenAI. Safety best practices for deploying AI agents. \url{https://platform.openai.com/docs/guides/safety-best-practices}, 2026.

\end{thebibliography}

\appendix

\section{Appendix: Implementation Details}

\subsection{Repository Structure}

\begin{lstlisting}
aegis/
  aegis-sidecar/         # Rust sidecar proxy (core)
  aegis-policy-engine/   # OPA/WASM policy engine
  aegis-audit-log/       # Immutable audit log
  aegis-gateway/         # Go-based multi-provider gateway
  aegis-control-plane/   # Centralized configuration (Go)
  aegis-cli/             # CLI tool
  aegis-common/          # Shared types, crypto, protobuf
  aegis-ebpf/            # eBPF programs + Rust loader
  deploy/                # K8s, Docker, Helm
  policies/              # OPA/Rego policy library
  integrations/          # Framework adapters
  tests/                 # E2E, security, chaos, benchmarks
  proto/                 # Protobuf definitions
  scripts/               # Build, benchmark, security scripts
  docs/                  # Architecture, security, protocol specs
\end{lstlisting}

\subsection{Architecture Diagram}

\begin{verbatim}
Agent (LangGraph/CrewAI/AutoGen)
        |
        v
+-------------------+
|  Governance        |
|  Sidecar (Rust)    |  -- Identity, Policy, Cost, Audit
|         |          |
|   +-----v-----+   |
|   |  eBPF     |   |  -- Non-bypassable infra-plane
|   |  Probes   |   |
|   +-----------+   |
+--------+----------+
         |
         v
+-------------------+
|  Agent Gateway    |  -- Multi-provider routing, rate limiting
|  (Go)             |
+--------+----------+
         |
     +---+----+
     v        v
  LLM APIs    Tools/Services
  (OpenAI,    (Databases, APIs,
   Anthropic)  Internal Services)
\end{verbatim}

\end{document}
